\documentclass[12pt]{article}
\usepackage[a4paper,margin=1in]{geometry}
\usepackage[T1]{fontenc}
\usepackage[utf8]{inputenc}
\usepackage{lmodern,microtype,amsmath,amssymb,graphicx,booktabs,array,longtable,placeins,chngcntr}
\usepackage{setspace,float}
\graphicspath{{../04_Figures/}}
\usepackage[style=apa,natbib=true,backend=biber,sortcites=true,sorting=nyt]{biblatex}
\addbibresource{../06_Bibliography/aeri_refs.bib}
\usepackage[colorlinks=true,citecolor=blue,linkcolor=blue,urlcolor=black]{hyperref}
\usepackage{caption}
\captionsetup{font=small,labelfont=bf}
\onehalfspacing
\setlength{\emergencystretch}{2em}
\newcommand{\clip}{\operatorname{clip}}
\newcommand{\ind}{\mathbf{1}}
\newcommand{\exhibit}[4]{\begin{figure}[htbp]\centering\includegraphics[width=\textwidth,height=.65\textheight,keepaspectratio]{#1}\caption{#2}\label{#3}\begin{singlespace}\small\textit{Notes:} #4\end{singlespace}\end{figure}}
\begin{document}
\begin{titlepage}\thispagestyle{plain}
\begin{center}
{\Large\bfseries Measuring Economic Risk in Africa: The African Economic Risk Index and African Risk Score\par}
\end{center}
\vspace{.35cm}
\begin{abstract}\singlespacing
We develop a multidimensional news-based framework for measuring economic risk across 54 African economies from 2015 to 2025. The African Economic Risk Index (AERI) uses graded article-level scores to capture the prevalence, severity and textual prominence of economic-risk reporting. Eight non-exclusive pillar indices identify where risk is concentrated, while the African Risk Score (ARS) combines bounded measures of intensity, breadth, severity and structural spread on a 0--100 scale. The analysis uses 21.30 million retained article records and 6,488 country-months meeting the principal descriptive coverage threshold. Exact-text deduplication removes 497,091 repeated records (2.33 percent) from the reconstructed frozen vintage; original and duplicate-adjusted AERI remain highly correlated (Pearson 0.9296; rank 0.9883). Aggregate AERI peaks in April 2020, while country histories show substantial heterogeneity in timing, persistence and composition. Alternative ARS weighting produces closely aligned scores. Macroeconomic content is assessed with pooled panel VARs and an external monetary-policy design. Inflation-risk readings precede higher headline inflation, exchange-rate depreciation is followed by higher AERI, and stronger real activity is followed by lower AERI. Externally identified U.S. monetary-policy tightening shocks generate larger subsequent AERI increases among economies with greater predetermined external-debt exposure. Coverage, weighting, exposure and falsification checks support the measurement framework. AERI, the pillars and ARS provide transparent, decomposable measures of reported economic risk for use alongside conventional macroeconomic statistics.
\end{abstract}
\noindent\small\textbf{Keywords:} Economic risk; Africa; newspapers; text as data; risk measurement; alternative data.\\
\textbf{JEL Classification:} E37, C55, C81, O55.
\end{titlepage}

\section{Introduction}\label{sec:introduction}
Official statistics remain the main source of information about inflation, output, fiscal conditions and external balances, but they do not capture all changes in the economic environment as they occur. News reports often record shortages, financing pressures, currency stress, price increases and production disruptions before the full effects appear in conventional data. We develop a way to measure that information systematically across African economies.

We construct the \textbf{African Economic Risk Index (AERI)} from article-level scores that combine economic relevance, severity and textual emphasis. For each country and month, AERI averages these scores over all retained articles and multiplies the result by 100. Non-relevant articles remain in the denominator, so the index rises when economic-risk reporting becomes more prevalent, more severe or more prominent relative to the overall news flow.

We also construct eight non-exclusive pillar indices covering inflation, foreign-exchange and external conditions, debt and fiscal conditions, the financial sector, growth and the real economy, commodities and supply, political economy, and climate and disasters. These indices show where measured risk is concentrated. A separate PCA summary records common movement across the pillars. Finally, the \textbf{African Risk Score (ARS)} combines AERI intensity, breadth of relevant reporting, average severity and the number of active domains on a 0--100 scale. AERI is the principal measure; the pillars and ARS provide additional information about the composition and breadth of the underlying news signal.

The data contain 21.30 million retained article records for 54 African economies from January 2015 through December 2025. We observe 7,127 country-months, of which 6,488 contain at least 100 retained articles and form the principal descriptive sample. Aggregate AERI reaches its sample maximum in April 2020. Country-level series, however, display considerable variation in both the timing and composition of risk. ARS is also insensitive to a transparent alternative weighting rule: the benchmark and PCA-weighted versions have a correlation of 0.995 in levels and 0.997 in ranks.

We evaluate the economic content of the measures in two ways. First, pooled panel vector autoregressions (PVARs) relate the news measures to inflation, exchange-rate movements and monthly GDP growth. Inflation-risk lags foreshadow higher inflation, depreciation is followed by higher AERI, and stronger real activity is followed by lower AERI. Second, we use high-frequency U.S. monetary-policy surprises from \citet{jarocinski2020} as an external source of variation and interact them with 2014 external-debt exposure. AERI rises more after a tightening in African economies that entered the sample with greater external exposure. The differential response survives alternative exposure measures, an alternative monetary-policy shock, leave-one-country-out estimation, and a robustness sample restricted to FOMC announcement months.

Our approach follows the newspaper-based measurement discipline of \citet{baker2016}, but the object of interest is broader economic risk rather than policy uncertainty. Three features distinguish the construction. Article scores are graded rather than binary; economic domains can overlap within an article; and ARS separates the intensity of the aggregate signal from its breadth, severity and structural spread. The same scoring rule is applied across 54 economies with very different media depth and economic structures, while coverage remains explicit in the reported data.

The paper proceeds as follows. Section~\ref{sec:literature} discusses related work on uncertainty, text-based indicators and high-frequency economic measurement. Section~\ref{sec:methodology} describes the construction of AERI, the pillar indices and ARS. Section~\ref{sec:results} reports their empirical properties and validation exercises. Section~\ref{sec:summary} concludes, and Section~\ref{sec:limitations} discusses the main limitations.

\section{Literature Review}\label{sec:literature}
\subsection{Economic Risk, Expectations, and Ex Ante Behaviour}
Economic decisions depend on prospective outcomes as well as current conditions. \citet{bloom2009} shows how uncertainty can delay investment and hiring, while \citet{jurado2015} measures uncertainty through the unpredictability of economic outcomes. These contributions motivate attention to information that appears before realised activity. A news-based risk index captures a different object: the salience and severity of economic concerns in reporting.

A clearly anticipated increase in financing costs may generate substantial risk reporting even when there is little uncertainty about whether it will occur. Conversely, a highly uncertain development may attract little coverage in a thin archive. AERI is built to capture the first of these margins: the amount and intensity of economic-risk content in the observed news flow. The empirical analysis then asks whether movements in that content bear systematic relationships to realised macroeconomic conditions.
\subsection{Media, Information Frictions, and Macroeconomic Beliefs}
Information reaches households and firms unevenly. \citet{mankiwreis2002} show how slow information diffusion can generate persistent inflation dynamics. \citet{shiller2017} emphasises narratives as a mechanism through which economic interpretations spread. News therefore provides a record of changing attention and beliefs, although it also responds to realised conditions. Examining both directions of the relationship between news-based risk and macroeconomic variables is important for interpretation.

A persistent inflation response following inflation-risk reporting is compatible with several mechanisms. News may record cost pressures before their full price effects, reflect expectations about those pressures, or respond to information also used by firms. The reverse direction is informative as well: if depreciation raises AERI, the index is registering realised stress as it enters the news environment. For that reason, the empirical design examines feedback from macroeconomic conditions alongside the forward-looking content of the news measures.
\subsection{Newspaper-based indices of uncertainty and risk}
\citet{baker2016} construct economic policy uncertainty from articles jointly referring to the economy, policy and uncertainty. \citet{caldara2022} develop a geopolitical-risk measure, and \citet{ahir2022} use uncertainty-related language in country reports to extend international coverage. Our framework retains the text-based approach but separates economic reporting into eight domains and incorporates severity and structural breadth. This provides information about the composition of economic stress beyond the incidence of a single class of uncertainty terms.

The common measurement lesson is that a textual event must be defined together with its denominator and economic interpretation. AERI uses a graded article score, whereas ARS combines bounded information on intensity, breadth, severity and structural spread. The eight-pillar structure then shifts the descriptive question from how much risk is reported to where it is concentrated. That distinction is useful when external financing, domestic prices and fiscal pressures move together but remain economically distinct.

AERI also differs substantively from economic policy uncertainty. Its dictionary spans prices, production, finance and other economic domains, including developments whose occurrence may be relatively certain, and its article scores vary with severity and textual emphasis. EPU, by contrast, is built around joint references to the economy, policy and uncertainty. The two approaches therefore share a news-based measurement tradition but represent different constructs and scales. Our contribution is a broader economic-risk measure implemented consistently across an African panel.
\subsection{External Identification and Global Monetary Shocks}\label{sec:literature_external_id}
High-frequency monetary-policy identification provides a useful way to validate whether a news-based risk measure reacts to shocks determined outside the measurement system. \citet{jarocinski2020} separate pure monetary-policy surprises from central-bank information shocks using the high-frequency co-movement of interest rates and equity prices around policy announcements. We use their pure monetary-policy component as an external shock and interact it with predetermined country exposure. The exercise asks a deliberately narrow question: does an independently identified global tightening produce larger changes in AERI where pre-existing external financial exposure implies greater vulnerability?

\subsection{Text-as-Data Methods and Interpretability Trade-offs}
\citet{gentzkow2019} treat textual measurement as an empirical design problem in which preprocessing, feature definition and interpretation affect the resulting economic variable. Dictionary-based scoring allows the contribution of an article to be traced to explicit rules. PCA then reduces the dimensionality of the aggregate inputs \citep{jolliffe2016}. It determines weights from covariance, rather than choosing them to maximise correlation with inflation or output. Economic interpretation continues to depend on the definitions of the inputs.

PCA enters only at the aggregation stage. Economic relevance is defined beforehand through the article-scoring rules. Standardisation prevents the numerical scale of an input from determining its principal-component loading, and converting the loadings back to the bounded input scale explains why a component with a smaller standardised loading can receive a larger final weight. We report the loadings, scaling bounds and final weights so that the empirical weighting step can be reproduced directly.
\subsection{Cross-country measurement challenges and African contexts}
Cross-country indices are sensitive to publication volume, source composition, language and archive access. Scaling by total article counts removes the direct level effect of media volume, while coverage diagnostics track variation in the underlying observation base. We also distinguish missing source months from observed months with zero measured risk. These choices are especially important when one statistical procedure is applied across heterogeneous media environments.

Comparability has both a time-series and a cross-sectional dimension. Within a country, changes in the available outlet set can affect the observed news mix; across countries, similar scores can arise from media systems with different degrees of economic-news specialisation. The common scoring rule provides definitional consistency and the minimum-coverage restriction limits the influence of very small denominators. The secondary empirical exercises further absorb persistent country differences and shared calendar shocks. Time-varying changes in national source coverage remain a separate measurement concern and are examined through the coverage diagnostics.
\subsection{African Economies and Alternative Sources of High-Frequency Information}
A growing literature uses alternative data to supplement conventional economic statistics where coverage is sparse or publication lags are long. \citet{yeh2020} use publicly available satellite imagery and deep learning to measure local economic well-being across Africa. \citet{buell2021} combine high-frequency indicators, including search and payments data, with statistical models to track economic activity in Sub-Saharan Africa. More recently, \citet{fotopoulou2026} examine machine-learning and satellite-data methods for nowcasting real GDP in data-constrained settings. These studies illustrate a broader measurement problem: economically relevant information often arrives at frequencies and in forms that are not captured by standard national accounts releases.

News data address another margin of the high-frequency measurement problem by recording the content and salience of economic concerns as they appear in public reporting. The eight-pillar structure follows price pressures, fiscal concerns, external financing conditions and other domains separately, including periods when the aggregate index changes little. We use AERI and ARS alongside official macroeconomic statistics and evaluate whether their movements line up with macroeconomic conditions in economically sensible directions.
\section{Measuring AERI and ARS}\label{sec:methodology}
We proceed in three stages: data collection, article-level scoring and country-period aggregation. From these steps we construct four related objects: AERI, eight overlapping pillar indices, an eight-pillar PCA common-factor composite, and the bounded ARS. Table~\ref{tab:measure_map} summarises their roles. AERI tracks scored economic-risk content relative to the news flow; the pillars locate that content across eight domains; the PCA composite gauges common movement across those domains; and ARS condenses intensity, breadth, severity and structural spread into a monitoring score.
\begin{table}[htbp]\centering\small\singlespacing
\caption{Measurement objects and their roles}\label{tab:measure_map}
\setlength{\tabcolsep}{4pt}
\begin{tabular}{p{2.5cm}p{4.7cm}p{6.1cm}}\toprule
Measure & Construction & Interpretation\\\midrule
AERI & $100$ times the sum of article-level risk scores divided by all retained articles in a country-month & Overall intensity of scored economic-risk content relative to the observed news flow; the principal economic-risk index\\
Pillar index & Same article-level risk score as AERI, with the numerator restricted to articles assigned to one of eight economic domains & Domain-specific diagnostic showing where measured economic risk is occurring; overlapping assignments preserve cross-domain connections\\
Eight-pillar PCA composite & Log-transformed, percentile-scaled pillar indices combined with fixed first-component weights & Common-factor diagnostic showing how strongly the eight risk domains move together; useful for distinguishing broad cross-domain stress from a concentrated pillar episode\\
ARS & Weighted average of bounded AERI intensity, relevant-article breadth, conditional severity and structural spread, scaled to 0--100 & Multidimensional monitoring score summarising how intense, widespread, severe and broad-based the measured risk is\\
\bottomrule\end{tabular}
\par\medskip\raggedright\footnotesize\textit{Notes:} AERI is the principal article-score index. The pillar indices locate the economic domains in which that risk appears. The PCA composite provides a companion diagnostic of common cross-domain movement, while ARS is a separate bounded monitoring score. Each measure is interpreted according to its stated construction and role in the framework.
\end{table}

\subsection{Data Collection}\label{sec:data_collection}
We assemble monthly country-level records from national news sources and organise retrieval by country, site and month, with bounded retrieval attempts and site-level diagnostics. This procedure makes the collection process repeatable, although coverage can still vary across outlets and repeated stories may remain. We therefore report corpus size in retained article records. Section~\ref{sec:deduplication} evaluates the sensitivity of AERI and ARS to repeated identical article bodies using a separately reconstructed, internally matched archive.

Our analytical vintage contains 21,303,573 article records across 54 African countries from January 2015 to December 2025. Of these, 5,896,776 meet the economic relevance criterion. The 7,127 country-month records include 639 months with fewer than 100 articles. To limit sensitivity to very small news denominators, the principal descriptive comparisons use the remaining 6,488 country-months with at least 100 retained articles.\footnote{Months below the threshold remain in the underlying series and coverage files. The threshold governs the principal cross-country descriptive comparisons and the ARS calibration sample, while preserving the lower-volume observations in the documented archive.} The median country-month contains 1,039 articles. Months for which the source archive is absent or unreadable are recorded as missing.

The country-month is the unit of aggregation. A widely syndicated event can therefore generate several article records, sometimes across multiple outlets. We keep record counts distinct from independent stories and apply the coverage threshold at the country-month level before making cross-country descriptive comparisons. Calibration, descriptive analysis and macroeconomic estimation use separately documented samples because their data requirements differ. Section~\ref{sec:deduplication} quantifies how much the published series changes when identical article bodies are collapsed within country-months.
\exhibit{Figure_1.pdf}{News coverage and eligible country observations}{fig:coverage}{The left panel reports total retained article records by month. The right panel counts countries with at least 100 articles. Article counts refer to retained records before story-level exact-text deduplication. Changes in the archive can affect measured coverage.}
\subsection{Data Pre-processing}\label{sec:preprocessing}
Preprocessing converts article text into economic relevance, severity, textual emphasis and pillar assignments. The same rules apply across the retained corpus. An article can activate several pillars because economic developments often span more than one domain.
\subsubsection{Text Normalisation}\label{sec:text_norm}
The headline and body form a single textual unit, $x_i$. Normalisation removes formatting differences before dictionary matching. Scoring is rule-based and uses observable textual features from this normalised text.
\subsubsection{Economic Relevance Filtering}\label{sec:relevance_filter}
Relevance is defined by
\begin{equation}R_i=\ind\{x_i\text{ matches the economic-relevance dictionary}\}.\label{eq:relevance}\end{equation}
Non-relevant articles receive zero economic-risk score but remain in the denominator. Relevant reporting also includes discussions of economic conditions expressed without acute-crisis language.

All retained articles enter the denominator, so the share of the news flow devoted to economic risk remains part of the signal. The separate breadth and severity diagnostics then recover, respectively, the prevalence of relevant reporting and its conditional intensity. For an observed source month with retained records, a zero score is therefore an empirical outcome; missing source months are handled separately.
\subsubsection{Severity Tiering}\label{sec:severity_tiering}
Let $\mathcal K_i\subseteq\{1,2,3,4\}$ contain the severity tiers matched in the text. We set
\begin{equation} S_i=\begin{cases}0,&R_i=0,\\ \max(\{1\}\cup\mathcal K_i),&R_i=1.\end{cases}\label{eq:severity}\end{equation}
A relevant article receives a baseline of one, with stronger matched language increasing the score. Where several cues occur, the highest tier applies. The resulting severity tier is an ordinal description of language strength.
\subsubsection{Textual Emphasis Weighting}\label{sec:confidence_weighting}
The textual weight combines a headline cue $T_i$, a body-length indicator $L_i$ for at least 1,200 normalised characters, and an indicator $M_i$ for at least two distinct active pillars:
\begin{equation}W_i=\clip(1+0.5T_i+0.2L_i+0.3M_i,1,2).\label{eq:confidence}\end{equation}
Here $\clip(x,a,b)=\min(\max(x,a),b)$. The weight is a bounded textual-emphasis adjustment based on placement, length and cross-domain content. Its coefficients belong to the article-scoring rule and are distinct from the PCA aggregation weights used later in the construction.
\subsubsection{Article-Level Risk Scoring}\label{sec:article_scoring}
The score is
\begin{equation} A_i=R_iS_iW_i,\qquad 0\leq A_i\leq8.\label{eq:article}\end{equation}
Each activated pillar contributes once to an article score; repeated phrases within a story therefore leave that pillar's contribution unchanged. Table~\ref{tab:keyword_example} shows representative terms from the production dictionaries and a worked example of how the matching rules map one article into AERI, pillar and ARS inputs.
\begin{table}[H]\centering\small\singlespacing
\caption{Illustrative keyword taxonomy used in article scoring}\label{tab:keyword_example}
\setlength{\tabcolsep}{4pt}
\begin{tabular}{p{2.6cm}p{5.7cm}p{6.1cm}}\toprule
Scoring layer & Illustrative terms from the production dictionary & Role in the computation\\\midrule
Economic relevance & economy, GDP, inflation, budget, debt, central bank, currency, bank, unemployment; French examples include \emph{économie}, \emph{PIB}, \emph{dette}, \emph{taux de change}; Portuguese examples include \emph{economia}, \emph{PIB}, \emph{dívida}, \emph{câmbio} & At least one match sets $R_i=1$. Non-matching articles receive a zero article score but remain in the country-month denominator.\\
Risk pillars & Inflation: inflation, cost of living, shortage; FX/external: exchange rate, currency, devaluation, reserves; debt/fiscal: debt, default, restructuring, budget deficit; growth: recession, slowdown, unemployment; analogous dictionaries define the other four pillars & Pillar matches identify the domain composition of a relevant article. Pillars are non-exclusive and repeated mentions do not increase the pillar weight.\\
Severity tiers & Tier 1: concern, risk, uncertainty; Tier 2: surge, spike, worsening, deterioration; Tier 3: crisis, collapse, emergency; Tier 4: default, hyperinflation, bank run, economic collapse & A relevant article receives baseline severity 1; when several severity cues occur, $S_i$ equals the highest matched tier.\\
Title/emphasis cues & crisis, default, devaluation, inflation, shortage, collapse, debt, restructuring, recession, bank, emergency & A title match adds 0.5 to $W_i$; a body of at least 1,200 normalised characters adds 0.2; two or more active pillars add 0.3. The weight is clipped to $[1,2]$.\\
\bottomrule\end{tabular}
\begin{singlespace}\footnotesize\raggedright
\textit{Notes:} The table reports representative terms rather than the full dictionary. The supplied production configuration contains explicit English, French and Portuguese terms. Terms are assigned to relevance, pillar, severity and title-hint dictionaries before aggregation, independently of the subsequent macroeconomic validation outcomes.
\end{singlespace}
\end{table}

\paragraph{Worked scoring example.}
Consider a hypothetical article titled ``Inflation surge after currency devaluation.'' Suppose the normalised body exceeds 1,200 characters and contains no severity cue above tier 2. The article matches the economic-relevance dictionary, the inflation and FX/external pillars, the tier-2 term ``surge,'' and a title-risk cue. Hence $R_i=1$, $S_i=2$, $T_i=L_i=M_i=1$, so
\[
W_i=\clip(1+0.5+0.2+0.3,1,2)=2,
\qquad A_i=R_iS_iW_i=4.
\]
If a country-month contains 100 retained articles and this is the only relevant article, this one record contributes $100(4/100)=4$ points to AERI. The inflation and FX/external pillar indices each receive 4 points because the article contributes once to each active pillar. The same hypothetical month has breadth $B=1/100=0.01$, conditional severity $V/4=2/4=0.50$, and structural spread $K/8=2/8=0.25$. With the fixed PCA ARS weights $(0.2625,0.2722,0.3223,0.1430)$ and intensity $\mathrm{AERI}/200=0.02$, the corresponding illustrative score is
\[
\mathrm{ARS}=100\{0.2625(0.02)+0.2722(0.01)+0.3223(0.50)+0.1430(0.25)\}=20.49.
\]
The example is purely mechanical and illustrates article-level scoring; monthly scores are formed from the full retained country-month article set.

\subsection{Construction of the African Economic Risk Index}\label{sec:aeri}
For country $c$ and month $t$, let $\mathcal A_{ct}$ denote the retained articles and $N_{ct}=|\mathcal A_{ct}|$. AERI is
\begin{equation}\mathrm{AERI}_{ct}=\frac{100}{N_{ct}}\sum_{i\in\mathcal A_{ct}}A_i.\label{eq:aeri}\end{equation}
The index measures scored economic content per article. It can exceed 100, with a theoretical upper bound of 800. Annual AERI weights monthly values by article counts, so the annual numerator and denominator correspond to the same article universe.

The construction has a useful decomposition. Define $\overline{SW}_{ct}^{\,R}=\sum_iR_iS_iW_i/\sum_iR_i$ whenever relevant articles are present. Then
\begin{equation}
 \mathrm{AERI}_{ct}=100B_{ct}\overline{SW}_{ct}^{\,R}.\label{eq:aeri_decomposition}
\end{equation}
If there are no relevant articles, AERI is zero for a non-empty month. Equation~\eqref{eq:aeri_decomposition} shows why rising AERI can reflect a larger share of economic-risk reporting, more severe or emphatic relevant reporting, or both. This decomposition also makes clear that breadth and intensity are related inputs to ARS. PCA is well suited to this setting because it summarises their covariance rather than treating the dimensions as statistically independent.
\subsection{Pillar-Level Decomposition and Diagnostics}\label{sec:pillars_diagnostics}
Let $I_{pi}$ indicate that article $i$ activates pillar $p$. The pillar index is
\begin{equation}P_{pct}=\frac{100}{N_{ct}}\sum_{i\in\mathcal A_{ct}}A_iI_{pi}.\label{eq:pillar}\end{equation}
The eight pillars can overlap.\footnote{This treatment parallels the category-specific newspaper indexes in \citet{baker2016}, where a single article may enter more than one category.} Each pillar lies between zero and AERI, and the sum can exceed AERI when a story contributes to several domains. Breadth, severity and structural coverage are
\begin{equation}B_{ct}=\frac{\sum_iR_i}{N_{ct}},\quad V_{ct}=\frac{\sum_iS_i}{\max(1,\sum_iR_i)},\quad K_{ct}=\sum_{p=1}^{8}\ind\{P_{pct}>0\}.\label{eq:diagnostics}\end{equation}
The raw pillar indices locate risk across domains. To summarise their common cross-domain movement, define $\ell_{pct}=\log(1+P_{pct})$ and
\begin{equation}Q_{pct}=\clip\left(\frac{\ell_{pct}-a_p}{b_p-a_p},0,1\right),\quad C_{pct}=100u_pQ_{pct},\quad H_{ct}=\sum_{p=1}^{8}C_{pct}.\label{eq:pca_pillars}\end{equation}
The bounds $a_p$ and $b_p$ are the first and ninety-ninth calibration percentiles, and $u_p$ are fixed weights obtained from the first principal component of the standardised transformed pillars. The resulting $H_{ct}$ is the eight-pillar PCA composite, which we use as a common-factor diagnostic. AERI can be high when overall scored risk content is elevated but concentrated in a few domains; the PCA composite rises most strongly when several transformed pillar indices move together. This distinction helps separate broad cross-domain stress from episodes driven mainly by one or two pillars.

The PCA composite adds information to the ARS structure term. The ARS component $K_{ct}/8$ counts active pillars; $H_{ct}$ also incorporates their transformed intensity and empirical common-factor weights. Table~\ref{tab:pillar_weights_full} reports the corresponding loadings, scaling bounds and weights.

\subsection{Construction of the African Risk Score}\label{sec:ars}
ARS combines four bounded inputs:
\begin{equation}X_{ct}=\left(\clip\left(\frac{\mathrm{AERI}_{ct}}{200},0,1\right),B_{ct},\frac{V_{ct}}4,\frac{K_{ct}}8\right)',\qquad\mathrm{ARS}_{ct}=100\sum_{j=1}^{4}w_jX_{jct}.\label{eq:ars}\end{equation}
The divisor 200 sets the saturation point for the intensity component. The four entries have direct interpretations: $\clip(\mathrm{AERI}/200,0,1)$ is bounded AERI intensity, $B$ is the share of retained articles that are economically relevant, $V/4$ scales conditional severity to the unit interval, and $K/8$ records the share of the eight domains that are active. ARS is therefore a bounded 0--100 summary of the reported risk environment, rising as risk-related reporting becomes more intense, prevalent, severe or structurally widespread.

The four components are allowed to covary. Breadth is already one margin of AERI through equation~\eqref{eq:aeri_decomposition}, but it enters ARS separately by design. This allows two country-months with similar AERI to receive different ARS values when one reflects narrow, highly intense reporting and the other reflects broader, multi-domain stress. PCA summarises the covariance among these diagnostics and reduces the dependence of the final score on arbitrary hand-set weights.

Standardise the calibration inputs as $Z_j=(X_j-\mu_j)/s_j$, using population standard deviations. Let $v$ be the unit-length first principal-component loading vector, oriented toward higher component values. The weights on the bounded input scale are
\begin{equation}v=\arg\max_{\|b\|=1}b'\widehat\Sigma_Zb,\qquad w_j=\frac{v_j/s_j}{\sum_kv_k/s_k}.\label{eq:weights}\end{equation}
Dividing by $s_j$ converts standardised loadings to coefficients on the bounded input scale. With positive loadings, ARS is a rescaled first component and is weakly increasing in its bounded inputs. The same conversion is applied separately to the transformed pillar vector $Q$ to obtain the fixed common-factor weights $u_p$ used in equation~\eqref{eq:pca_pillars}. Calibration rejects near-constant inputs or materially mixed-sign loadings rather than converting negative loadings to absolute values.

We calibrate ARS using 6,238 country-months with at least 100 articles during January 2015--July 2025. ARS weights are 26.25 percent for intensity, 27.22 percent for breadth, 32.23 percent for severity and 14.30 percent for structure. PC1 explains 62.22 percent of standardised variation in the four ARS inputs and 52.50 percent for the transformed eight-pillar system. Weights remain fixed between recalibrations. Historical scores remain retrospective because calibration uses information through July 2025. In the eligible 2015--2025 sample, the median month activates seven of the eight pillars, while the median relevant-article share is 25.8 percent. These diagnostics help explain why ARS can move differently from AERI even though both are generated from the same underlying article corpus.

Tables~\ref{tab:component_weights_full} and \ref{tab:pillar_weights_full} report the two PCA calibrations. Severity receives the largest ARS weight even though intensity has the largest standardised loading because equation~\eqref{eq:weights} converts standardised loadings back to coefficients on the bounded input scale before normalisation. Table~\ref{tab:pillar_weights_full} documents how the transformed domain indices load on their common factor and makes the construction used in Figure~\ref{fig:weighting_sensitivity} auditable. Debt and fiscal risk and financial-sector risk receive the largest common-factor weights, at about 17.3 percent each, while commodity and supply risk receives 6.48 percent.

Calibration pools eligible country-month rows without imposing equal total country weights. Countries with more eligible months consequently contribute more observations. Holding the calibration fixed makes subsequent changes in a score attributable to changed inputs rather than to month-by-month re-estimation of its weights. A future recalibration should be treated as a new measurement vintage, with the associated parameters retained. Because ARS compresses several diagnostics into one score, the individual pillar and component series remain necessary for interpretation; information outside the common movement of the four ARS inputs may still be economically important.
\begin{table}[htbp]\centering\small\singlespacing
\caption{ARS component weights and standardised loadings}\label{tab:component_weights_full}
\setlength{\tabcolsep}{4pt}
\begin{tabular}{lrrr}\toprule
Component & Benchmark (\%) & PCA (\%) & PC1 loading\\\midrule
Intensity & 40.00 & 26.25 & 0.5803\\
Breadth & 25.00 & 27.22 & 0.5502\\
Severity & 20.00 & 32.23 & 0.3921\\
Structure & 15.00 & 14.30 & 0.4548\\
\bottomrule\end{tabular}
\par\medskip\raggedright\footnotesize\textit{Notes:} The 6,238 calibration observations cover January 2015--July 2025 and have at least 100 articles. PC1 explains 62.22 percent of standardised variation. Weights on the bounded input scale differ from standardised loadings; equation~\eqref{eq:weights} gives the conversion.
\end{table}
\begin{table}[htbp]\centering\small\singlespacing
\caption{PCA weights, scaling bounds and loadings for the eight pillars}\label{tab:pillar_weights_full}
\setlength{\tabcolsep}{4pt}
\begin{tabular}{lrrrr}\toprule
Pillar & Weight (\%) & Log lower & Log upper & PC1 loading\\\midrule
Inflation & 12.95 & 0.0000 & 2.9629 & 0.4052\\
FX and external & 13.21 & 0.0000 & 4.2958 & 0.3584\\
Debt and fiscal & 17.29 & 0.0000 & 4.3550 & 0.4071\\
Financial sector & 17.30 & 0.0000 & 4.5177 & 0.4219\\
Growth and real economy & 11.90 & 0.0000 & 2.6486 & 0.3861\\
Commodity and supply & 6.48 & 0.0000 & 1.0248 & 0.1289\\
Political economy & 9.92 & 0.0000 & 2.9873 & 0.3063\\
Climate and disaster & 10.95 & 0.0000 & 2.3585 & 0.3227\\
\bottomrule\end{tabular}
\par\medskip\raggedright\footnotesize\textit{Notes:} Bounds are the first and ninety-ninth calibration percentiles of $\log(1+P_p)$. PC1 explains 52.50 percent of standardised variation. The weights quantify each transformed pillar's contribution to the common cross-domain factor after scale adjustment; economic interpretation remains anchored in the domain-specific series and results. Calculations use unrounded weights, which sum to one.
\end{table}
\section{Empirical Properties and Evaluation}\label{sec:results}
\subsection{Aggregate Risk Dynamics}\label{sec:aggregate_risk}
Figure~\ref{fig:trends} presents the two principal measures used throughout the paper. The left panel shows AERI as defined in equation~\eqref{eq:aeri}; the right panel shows ARS from equation~\eqref{eq:ars}, using the fixed PCA weights in Table~\ref{tab:component_weights_full}. For each month, the figure plots the unweighted cross-country median with interquartile ranges as shaded bands. The eight-pillar PCA composite is analysed separately in Figure~\ref{fig:weighting_sensitivity}.

The monthly median AERI reaches its sample maximum of 68.94 in April 2020, coinciding with the acute phase of the COVID-19 pandemic and the associated economic disruption. The shaded bands show dispersion across eligible countries. Because the set of eligible countries changes with archive coverage, aggregate movements are read alongside Figure~\ref{fig:coverage}.

Across 6,488 eligible country-months, AERI has a mean of 47.89 and a standard deviation of 27.56. ARS has a mean of 36.59 and a standard deviation of 9.66. Their different scales follow directly from construction: AERI preserves article-score intensity, while ARS combines bounded components. The two measures are read jointly, with AERI describing the level of the scored signal and ARS its broader configuration.
\exhibit{Figure_2.pdf}{Aggregate AERI and ARS dynamics}{fig:trends}{Lines are unweighted country medians; shaded regions are interquartile ranges. The annotation marks the COVID-19 pandemic shock at the maximum monthly median AERI in the sample (April 2020, 68.94). Only country-months with at least 100 articles enter. Neither series is weighted by GDP. ARS uses the fixed PCA calibration.}

Figure~\ref{fig:weighting_sensitivity} provides two complementary robustness checks. Panel A compares AERI with the eight-pillar PCA composite, which summarises common movement across the transformed domain indices. Their monthly-median correlation is 0.558 over 2015--2025, and both reach their maximum in April 2020. The common peak indicates that the dominant AERI episode coincided with unusually broad movement across risk domains. Outside that episode the series diverge in informative ways. May 2022, for example, is the second-highest monthly median of the pillar composite but falls outside the five highest AERI months, pointing to stronger cross-domain co-movement than aggregate article-score intensity alone would suggest.

Panel B isolates the weighting choice. The benchmark 40/25/20/15 score and the PCA-weighted score use the same inputs and observations; their monthly cross-country medians correlate 0.975 and both attain their 2015--2025 maximum in April 2020. The broad ARS chronology is therefore stable under a substantially different weighting scheme.
\exhibit{Figure_3.pdf}{Weighting sensitivity of the AERI-related and ARS measures}{fig:weighting_sensitivity}{Panel A compares AERI with the eight-pillar PCA common-factor composite. Panel B compares benchmark ARS, using 40/25/20/15 weights on intensity, breadth, severity and structure, with PCA-weighted ARS. Lines are monthly unweighted country medians for country-months with at least 100 retained articles. The vertical reference marks April 2020. The monthly-median correlations are 0.558 for AERI versus the eight-pillar PCA composite and 0.975 for benchmark versus PCA-weighted ARS.}

\begin{table}[htbp]\centering\small\singlespacing
\caption{Distribution of AERI and alternative ARS specifications}\label{tab:comparison}
\setlength{\tabcolsep}{4pt}
\begin{tabular}{lrrrrr}\toprule
Measure & Mean & SD & Median & Minimum & Maximum\\\midrule
AERI & 47.89 & 27.56 & 46.91 & 0.00 & 300.16\\
ARS (PCA weights) & 36.59 & 9.66 & 37.27 & 0.00 & 85.37\\
Benchmark ARS & 35.42 & 10.71 & 36.01 & 0.00 & 88.93\\
\bottomrule\end{tabular}
\par\medskip\raggedright\footnotesize\textit{Notes:} Common sample: 6,488 country-months with at least 100 articles, 2015--2025. The benchmark and PCA-weighted ARS have Pearson correlation 0.9950, rank correlation 0.9971, and mean absolute difference 1.49 points. The comparison changes only the ARS weights; AERI itself has no PCA weighting stage.
\end{table}

Table~\ref{tab:comparison} reinforces the visual evidence in Figure~\ref{fig:weighting_sensitivity}. The PCA-weighted and benchmark ARS have very similar central tendencies and nearly identical rankings, although the PCA specification is modestly less dispersed (standard deviation 9.66 versus 10.71). Changing the weights therefore shifts the scale slightly while leaving the relative ordering of country-months largely intact. AERI is considerably more dispersed because it preserves the original article-score intensity and remains on its own scale.
\subsection{Annual Africa Economic Risk Indices by Country (2015--2025)}\label{sec:aeri_heatmap}
Figure~\ref{fig:annual} follows annual AERI for the ten economies examined in the country discussion. The common colour scale makes differences in intensity and persistence visible without imposing a common economic explanation on each episode. Annual values aggregate the monthly article-score numerator and denominator.
Among countries with at least nine eligible months in 2025, the highest annual AERI values are Zambia (112.40), Liberia (99.59), Lesotho (91.14). Annual values are article-count-weighted averages of monthly AERI scores. The resulting ordering describes the intensity of observed economic-risk reporting across the retained news archive.
We complement the annual comparison with exact monthly summary statistics for the same ten-country group in Appendix Table~\ref{tab:country_aeri}; the full all-country summary and monthly country paths are reported in the accompanying Supplementary Material. These outputs document differences in levels, variability and the timing of peaks that the aggregate median and annual averages conceal. Within the ten-country group, Nigeria has the highest mean monthly AERI, while Ethiopia exhibits the greatest time-series dispersion. The country evidence is descriptive and uses the same AERI definition as the aggregate analysis.
\exhibit{Figure_4.pdf}{Annual AERI in selected African economies}{fig:annual}{The figure reports annual AERI for the ten illustrative economies over 2015--2025. Monthly AERI values are aggregated using monthly retained-article counts, so the annual statistic preserves the article-score numerator and denominator. The ten countries form the illustrative country group used in the descriptive discussion, and the underlying annual file covers all 54 economies. The common colour scale highlights substantial cross-country differences in levels, persistence and the timing of peaks.}

Figure~\ref{fig:annual} shows substantial heterogeneity beneath the aggregate African pattern. Several economies, including Angola, Egypt, Ethiopia, Kenya and South Africa, record clear increases in annual AERI in 2020 relative to their own pre-pandemic averages, consistent with the common shock visible in Figure~\ref{fig:trends}. The response is nevertheless far from uniform. Nigeria remains at a comparatively high level throughout much of the sample and reaches its annual maximum later, in 2024, while Morocco displays a gradual post-2020 rise rather than a single pandemic-centred jump. \mbox{Ethiopia peaks in 2019 and remains elevated in 2020.} These differences are important for interpretation: the continental median captures common timing, whereas the country series reveal persistence, delayed adjustment and country-specific episodes that are averaged away in the aggregate.

\subsection{Pillar-Level Decomposition}\label{sec:pillar_results}
The pillar indices decompose the aggregate signal by economic domain. Figure~\ref{fig:pillars} reports the pooled mean of each raw pillar index across the 6,488 eligible country-months, and Table~\ref{tab:pillar_profile} shows how the domain profile changes across broad subperiods. Because an article may activate several domains, these are overlapping measures of reported risk content rather than shares of a fixed total.
\exhibit{Figure_5.pdf}{Average domain-specific pillar indices}{fig:pillars}{Bars are ordered by magnitude and report pooled means of the eight raw pillar indices across 6,488 country-months with at least 100 retained articles. Each pillar uses the same graded article score as AERI, with the numerator restricted to articles assigned to that domain. Pillars are non-exclusive, so a single article may contribute to several domain indices. Financial-sector and debt/fiscal reporting have the largest pooled means in this vintage.}

Financial-sector risk has the largest pooled mean, followed by debt and fiscal risk. FX/external, political-economy and inflation risks form a middle group; growth and real-economy risk sits below them, while climate/disaster and commodity/supply risks are the least prominent on average. The ordering reflects both the incidence of economic concerns and the way those concerns enter the news. Raw dictionary size offers little explanation for the financial-sector ranking: that pillar contains nine top-level terms or phrases, the same number as six other pillars, while commodity/supply contains seven. Differences in term prevalence, newsworthiness and outlet specialisation are more plausible sources of cross-pillar level differences. We therefore emphasise within-pillar changes over time and use the levels primarily to describe the composition of reported risk.

\begin{table}[htbp]\centering\small\singlespacing
\caption{Pillar-specific economic-risk content over time}\label{tab:pillar_profile}
\setlength{\tabcolsep}{4pt}
\begin{tabular}{lrrrr}\toprule
Pillar & 2015--2019 & 2020 & 2021--2025 & Overall\\\midrule
Inflation risk & 3.47 & 3.49 & 5.33 & 4.34\\
FX and external risk & 5.59 & 5.11 & 4.48 & 5.03\\
Debt and fiscal risk & 9.47 & 9.64 & 7.65 & 8.64\\
Financial-sector risk & 11.75 & 12.87 & 10.90 & 11.46\\
Growth and real economy risk & 3.14 & 4.85 & 2.85 & 3.17\\
Commodity and supply risk & 0.07 & 0.03 & 0.13 & 0.09\\
Political-economy risk & 5.27 & 5.51 & 4.50 & 4.94\\
Climate and disaster risk & 1.25 & 1.11 & 1.27 & 1.25\\
\bottomrule\end{tabular}
\par\medskip\raggedright\footnotesize\textit{Notes:} Entries are means of the raw pillar indices on the same score-per-article scale used in AERI, computed over eligible country-months. The pillars are non-exclusive: one article may contribute to more than one domain, and the rows describe overlapping domain-specific content. The period split is descriptive and highlights changes in the composition of reported economic risk.
\end{table}

The time profile reinforces this interpretation. Financial-sector and debt/fiscal risk remain the two largest domains in each reported period, but the composition changes around major episodes. Growth and real-economy risk rises from 3.14 in 2015--2019 to 4.85 in 2020, while inflation risk rises from 3.47 before the pandemic to 5.33 in 2021--2025. FX/external and debt/fiscal risk are somewhat lower on average in 2021--2025 than before 2020. The pillar system therefore distinguishes changes in the composition of risk that would be obscured by the aggregate AERI series alone.
\subsection{From AERI to ARS: Intensity, Breadth, and Structure}\label{sec:aeri_vs_ars}
ARS separates the intensity of reporting from its prevalence and severity. Structural breadth measures the number of domains with positive scores, so it can remain high even when reporting intensity is low. The distinction helps identify whether an increase in AERI reflects a concentrated episode or wider economic coverage.

A fixed-weight benchmark using 40/25/20/15 percent for intensity, breadth, severity and structure has a pooled Pearson correlation of 0.995 with PCA-weighted ARS and a rank correlation of 0.9971. Their mean absolute difference is 1.49 points, showing that the broad ARS pattern is highly insensitive to the weighting choice. The pooled correlation between AERI and PCA-weighted ARS is lower, at 0.911. They clearly share a common signal, but ARS changes the scale and gives explicit weight to breadth, severity and structural spread.

Table~\ref{tab:comparison} provides the corresponding distributional comparison. We hold the sample and component definitions fixed, so the exercise isolates the effect of the weights themselves. The PCA specification supplies a reproducible empirical weighting criterion, while the 40/25/20/15 benchmark offers a transparent reference point. The two schemes diverge most when intensity and severity move differently, because those components receive the largest differences in weight. A country-month with high conditional severity but only moderate AERI intensity will tend to score higher under the PCA weights; an episode dominated by unusually high intensity with less exceptional severity will tend to score higher under the benchmark.
\subsection{Evaluation of AERI and ARS}\label{sec:validation}
We evaluate AERI and ARS in five ways, following the broad measurement logic in \citet{baker2016}. We verify the accounting identities and document news coverage; compare AERI with the eight-pillar common factor and ARS with a transparent alternative weighting rule; examine correspondence with salient economic episodes; estimate selected macroeconomic relationships in PVARs; and use externally identified U.S. monetary-policy shocks interacted with predetermined country exposure. These exercises examine reproducibility, aggregation sensitivity and economic content using distinct sources of variation.

\subsubsection{Accounting, coverage and aggregation robustness}
We reconcile AERI directly to its article-score numerator and total-article denominator. Equation~\eqref{eq:aeri_decomposition} then separates the relevance-share margin from conditional severity and emphasis. Pillar scores use the same article score with overlapping domain indicators, and ARS combines four bounded diagnostics under a frozen calibration. This sequence provides an auditable path from article text to the published country-month measures.

Coverage accounting is equally important. The 2015--2025 vintage contains 7,127 country-month records, of which 6,488 meet the 100-article threshold used for the principal descriptive comparisons. Lower-volume months remain in the documented archive but are omitted from those comparisons because ratio-based measures become unstable with very small denominators. Figure~\ref{fig:coverage} shows the resulting variation in the observation base. As an archive-depth diagnostic, we regress AERI on log retained article volume within the eligible sample. The within-country correlation is 0.035; the article-volume coefficient is 1.14 with country fixed effects ($p=0.702$) and $-1.47$ after adding year-month fixed effects ($p=0.640$). Appendix Table~\ref{tab:archive_depth} reports the specifications. The estimates show little association between AERI and archive size within countries. Changes in source composition and reporting practice remain separate concerns.

We assess ARS weight sensitivity by replacing the PCA weights with the benchmark 40/25/20/15 rule. The broad score pattern is almost unchanged: Pearson correlation is 0.9950, rank correlation is 0.9971 and the mean absolute difference is 1.49 points. With the sample and component definitions held fixed, the exercise attributes the remaining differences to the weighting choice.

\subsubsection{Syndicated stories and exact-text deduplication}\label{sec:deduplication}
Repeated publication can amplify a country-month signal when the same high-scoring text appears in several outlets. We collapse identical decoded article bodies within each country-month, retaining one representative per duplicate group.\footnote{Exact matching is applied to decoded bodies of at least 200 characters. Appearances in different countries or months remain separate news exposures, while minor edits, translations and rewritten versions remain distinct records.}

The correction is applied algebraically to the frozen analytical vintage. Each removed record is scored with the original dictionaries, and its contribution is subtracted from the archived numerator, relevance and severity counts and each pillar numerator. The all-article denominator is reduced by the same record count. Source files are read-only, and duplicate counts, representative examples and file fingerprints are retained in the replication extension.

Across 21,303,020 frozen-vintage records represented in 7,126 reconstructed country-months, the rule removes 497,091 repetitions (2.33\%), including 186,379 cross-outlet copies. The reconstruction is 553 records and one Togo country-month short of the archived aggregate because 42 supplied source files were empty or unreadable. For months with a partial count mismatch, unmatched frozen-vintage records remain in the denominator and only directly observed duplicates are removed. On the common coverage sample, original and duplicate-adjusted AERI have Pearson correlation 0.9296 and rank correlation 0.9883. The mean absolute change is 2.32 AERI points, and the 95th-percentile absolute change is 8.30. Fixed-weight ARS has correlation 0.9680 and mean absolute change 0.61 points. 69 country-months that originally met the 100-article rule fall below it after removal.

\input{../05_Tables/Table_7}

Country-level results show where duplication matters most. The median within-country AERI correlation is 0.9992, and 48 of 53 countries retain the same peak month on the common sample. Sierra Leone's original November 2016 peak falls from 269.72 to 110.77 after removal; Rwanda's August 2021 peak falls from 95.43 to 63.35; and Mauritius's May 2020 peak rises from 91.64 to 111.79 because duplicate removal changes both the numerator and the denominator. At the cross-country level, the median still peaks in April 2020. The exercise therefore supports the stability of the broad chronology while identifying a small number of country-specific spikes that are sensitive to identical reprints. The PVAR and external-shock specifications continue to use the original benchmark vintage.

\subsubsection{Event correspondence}
The aggregate AERI series reaches its sample maximum in April 2020, when the monthly cross-country median is 68.94. The timing coincides with the acute phase of the COVID-19 pandemic and provides a simple check on the economic content of the index. Country peaks occur in different months, however. Reading the aggregate series together with country paths and pillar composition reveals whether a spike reflects a broad continental episode or a more localised event.

\subsubsection{PVAR-based external-consistency evidence}
We use pooled PVARs and impulse-response functions as an external-consistency test of the measurement framework. The exercise centres on three directional implications: higher inflation-risk reporting should be followed by stronger realised inflation; an adverse exchange-rate innovation should raise measured economic risk; and stronger real activity should reduce it. Monthly GDP growth provides the measure of real activity in the PVAR system. Table~\ref{tab:selected_external_consistency} summarises the three relationships. All three joint lag tests are statistically significant at the paper's 10 percent criterion, with the inflation-risk and depreciation relationships also significant at 5 percent. At the reported horizons, the pointwise 90 percent intervals exclude zero in the predicted direction.\footnote{The bootstrap intervals are pointwise across horizons. They quantify horizon-specific uncertainty; simultaneous coverage of an entire impulse-response path would require a different band construction.}
\begin{table}[htbp]\centering\small\singlespacing
\caption{Directional PVAR external-consistency tests}\label{tab:selected_external_consistency}
\setlength{\tabcolsep}{4pt}
\begin{tabular}{@{}p{4.25cm}crp{6.0cm}@{}}\toprule
Relationship & Expected & Lag-test $p$ & Key response (90\% interval)\\\midrule
Inflation risk $\rightarrow$ headline inflation & $+$ & 0.0299 & M6: $+0.531$ pp $[0.287,\ 0.825]$\\
Depreciation $\rightarrow$ AERI & $+$ & 0.0142 & M2: $+0.0166$ SD $[0.0060,\ 0.0241]$\\
GDP growth $\rightarrow$ AERI & $-$ & 0.0624 & M1: $-0.181$ SD $[-0.338,\ -0.068]$\\
\bottomrule\end{tabular}
\par\medskip\raggedright\footnotesize\textit{Notes:} Expected signs are implied by the economic interpretation of the measures rather than chosen from the estimated responses. All three joint lag tests are statistically significant at the paper's 10 percent criterion, and the reported pointwise 90 percent intervals exclude zero in the expected direction. The inflation-risk and depreciation lag tests are also significant at 5 percent. Exact $p$-values are shown so readers can apply alternative significance conventions. These are conditional dynamic associations under the stated recursive ordering; structural identification is examined separately with the external monetary-policy shock design.
\end{table}


The strongest forward-looking result concerns the inflation-risk pillar. In the GDP-augmented PVAR, the three inflation-risk lags are jointly significant in the headline-inflation equation ($p=0.0299$). A one-within-country-standard-deviation inflation-risk innovation raises headline inflation by 0.111 percentage points on impact, 0.288 after one month and 0.392 after two months. The response continues to build and reaches 0.531 percentage points in month six, with a 90 percent interval of $[0.287,0.825]$. The pointwise 90 percent band stays above zero at every horizon from impact through month twelve. The hump-shaped profile indicates that the inflation-risk pillar contains information about sustained price pressure beyond the contemporaneous inflation reading. In the terminology of \citet{baker2016}, the news-based innovation \emph{foreshadows} subsequent inflation.
\exhibit{Figure_6.pdf}{Headline inflation response to an inflation-risk innovation}{fig:inflation_irfs}{The figure uses the 36-country GDP-augmented PVAR(3), estimated on 4,286 observations. The innovation is a one-within-country-standard-deviation increase in the article-score inflation-risk pillar. Bands are pointwise 90 percent intervals from 400 country-cluster bootstrap draws. The three inflation-risk lags are jointly significant in the headline-inflation equation ($p=0.0299$); the response is positive at every reported horizon and peaks at 0.531 percentage points in month 6.}

Figure~\ref{fig:aeri_macro_feedback} turns the validation question around and asks whether realised macroeconomic conditions feed into AERI in economically sensible directions. The exchange-rate result is particularly clear. A one-percentage-point innovation in monthly depreciation raises AERI by 0.0098 within-country standard deviations on impact, with a 90 percent interval of $[0.0033,0.0153]$. After a small month-one response, AERI rises to 0.0166 standard deviations in month two, with a 90 percent interval of $[0.0060,0.0241]$. The three depreciation lags are jointly significant in the AERI equation ($p=0.0142$). The response is short-lived rather than permanent, but its sign, timing and joint significance are consistent with currency deterioration broadening the economic-risk content of the news as exchange-rate pressure passes into import costs, financing conditions and expectations.

The real-activity response provides a complementary sign test. A one-percentage-point innovation in monthly GDP growth has a small and imprecise impact response, but AERI falls sharply one month later by 0.181 within-country standard deviations, with a 90 percent interval of $[-0.338,-0.068]$. Point estimates remain negative thereafter and gradually return toward zero. The joint three-lag Wald test is $p=0.0624$, which is statistically significant at the 10 percent level. The horizon-specific response and the joint lag test therefore point in the same direction: stronger activity is followed by less intense economic-risk reporting once information about improved real conditions is absorbed into the news flow.

The three IRFs point in the expected directions: inflation-risk reporting precedes higher realised inflation, currency depreciation is followed by higher AERI, and stronger real activity is followed by lower AERI. The forward inflation system places the risk measure first, while the feedback system places inflation, depreciation and GDP growth first. The sign pattern is therefore visible under both recursive orderings. We interpret these estimates as conditional dynamic relationships; structural identification comes from the separate external-shock exercise below.
\exhibit{Figure_7.pdf}{AERI responses to exchange-rate and real-activity innovations}{fig:aeri_macro_feedback}{Reverse 36-country GDP-augmented PVAR(3), estimated on 4,286 observations. Depreciation raises AERI, peaking at 0.0166 within-country standard deviations in month 2; GDP growth is followed by a $-0.181$ standard-deviation AERI response in month 1. Shading denotes pointwise 90 percent country-cluster bootstrap intervals. The opposite signs are consistent with higher reported risk under currency stress and lower reported risk after stronger activity.}

The PVAR evidence is intended to establish economic consistency rather than to catalogue every transmission channel. The paper's main contribution remains the construction and interpretation of AERI, the eight pillars and ARS.

\clearpage
\subsubsection{Externally identified monetary-policy exposure test}\label{sec:external_identification}
To complement the recursive PVAR, we turn to a shock measured outside the AERI system. Following \citet{jarocinski2020}, we use the pure U.S. monetary-policy surprise extracted from high-frequency movements in interest rates and stock prices around Federal Reserve announcements. We aggregate the surprise to the month and standardise it over 2015--2025. Country exposure is fixed at 2014 external debt stocks as a share of GNI from the World Bank International Debt Statistics \citep{worldbankdebt}, so the exposure measure predates the AERI sample.

For country $i$ and horizon $h$, the preferred local projection is
\begin{equation}
\begin{aligned}
\mathrm{AERI}_{i,t+h}-\mathrm{AERI}_{i,t-1}
={}&\alpha_i^h+\gamma_t^h+\beta_h\left(MP_t\times E_i^{2014}\right)\\
&+\sum_{\ell=1}^{3}\theta_{h\ell}\Delta\mathrm{AERI}_{i,t-\ell}
+\sum_{\ell=1}^{3}\psi_{h\ell}\left(MP_{t-\ell}\times E_i^{2014}\right)
+\varepsilon_{i,t+h}.
\end{aligned}
\label{eq:external_id}
\end{equation}
where $MP_t$ is a one-standard-deviation pure monetary-policy tightening and $E_i^{2014}$ is standardised $\log(1+\text{external debt/GNI})$. Country fixed effects absorb persistent differences in average risk reporting, while year-month fixed effects absorb the common global shock, COVID-19 and other disturbances shared across countries in a given month. The identifying variation is therefore cross-country exposure to the same externally measured tightening. Standard errors are two-way clustered by country and month following \citet{cameron2011}. We use finite-cluster reference inference and report restricted wild-cluster-bootstrap results as an additional finite-sample check.\footnote{The identifying cross-section contains 29 country clusters, so the finite-cluster $t$ and $F$ references use 28 denominator degrees of freedom. The restricted wild-cluster bootstrap follows \citet{mackinnon2017}.}

Our primary sample contains 29 African economies with both AERI and 2014 external-debt exposure. Figure~\ref{fig:external_identification} shows a delayed differential response. Four months after a one-standard-deviation tightening, a country one standard deviation more exposed before the sample records a 0.272-point larger cumulative AERI increase. The finite-cluster 90 percent interval is $[0.108,0.436]$ and $p=0.0086$. By month nine the differential reaches 0.384 points, with interval $[0.085,0.684]$ and $p=0.0373$. Restricted wild-cluster-bootstrap $p$-values are 0.042 and 0.034 at months four and nine, respectively. The response is delayed rather than contemporaneous, a pattern consistent with a gradual adjustment of the reported risk environment following tighter external financing conditions.

\exhibit{Figure_8.pdf}{AERI response to an externally identified U.S. monetary-policy tightening by predetermined external-debt exposure}{fig:external_identification}{The sample contains 29 African economies with AERI and 2014 external-debt exposure. The line reports $\hat\beta_h$ from equation~\eqref{eq:external_id}; shaded bands are pointwise 90 percent finite-cluster intervals using 28 denominator degrees of freedom. The shock is the standardised pure monetary-policy surprise from \citet{jarocinski2020}, and exposure is standardised $\log(1+\text{2014 external debt/GNI})$. The dependent variable is the cumulative change in AERI from month $t-1$ through horizon $h$. Country and year-month fixed effects and three lags of the outcome and interaction are included; standard errors are two-way clustered by country and month. The differential response is positive and statistically significant at months 4 and 9 under the finite-cluster reference.}

\input{../05_Tables/Table_9}

The exposure response is also stable across reasonable changes in exposure, sample composition and shock definition; Table~\ref{tab:external_identification} reports the individual specifications. Every month-four and month-nine estimate remains positive and significant at the paper's 10 percent criterion under finite-cluster inference. The month-four coefficient also remains significant at 5 percent across all six reported variants. The month-nine evidence is somewhat more sensitive in magnitude, while retaining a positive sign and 10 percent significance across all six variants. Restricting the sample to FOMC announcement months leaves the result intact.\footnote{This robustness specification changes the time observations only; the cross-section remains the same set of African economies.} Leave-one-country-out estimates tell the same story. At month four the coefficient ranges from 0.214 to 0.308, with all 29 estimates significant at 10 percent and 28 at 5 percent. At month nine it ranges from 0.282 to 0.455; all 29 are significant at 10 percent and 24 at 5 percent.

We next examine falsification tests. Current monetary-policy surprises have little relationship with AERI changes dated one to six months earlier; finite-cluster $p$-values range from 0.156 to 0.807. The central-bank-information component of the same announcements produces no comparable debt-exposure pattern at months four and nine ($p=0.153$ and $p=0.125$). The joint future-shock interaction placebo yields $F(6,28)=1.884$ with $p=0.119$, and the restricted wild-cluster bootstrap gives $p=0.221$.\footnote{The corresponding large-sample Wald statistic is $\chi^2(6)=11.30$ with asymptotic $p=0.079$. With 29 country clusters, the finite-cluster and restricted wild-cluster-bootstrap references provide the more appropriate finite-sample benchmarks.} Taken together, the placebos support the interpretation that the observed differential response is linked to the externally measured tightening. The exposure interaction should still be read with care. The 2014 debt ratio is predetermined and observational, and it may be correlated with exchange-rate regimes, trade openness, institutional quality and other persistent country characteristics. The coefficient therefore measures how the response to the external shock varies with pre-existing debt exposure. Establishing a causal debt-transmission channel would require a design that isolates exposure from these persistent country characteristics.

\FloatBarrier

\section{Conclusion}\label{sec:summary}
We develop two news-based measures of economic risk for 54 African economies over 2015--2025 and use them to study the level, composition and macroeconomic content of reported risk. AERI aggregates graded article-level scores for economic relevance, severity and textual emphasis; the eight pillar indices preserve information about the domains in which risk is concentrated; and ARS summarises intensity, breadth, severity and structural spread on a bounded scale. Applied to more than 21 million retained article records, the measures reveal substantial common variation alongside pronounced differences in the timing and composition of country-level risk. The macroeconomic exercises point in the same direction: inflation-risk readings precede movements in headline inflation, exchange-rate depreciation is followed by higher AERI, stronger real activity is followed by lower AERI, and an externally identified U.S. monetary-policy tightening is followed by a larger AERI increase in African economies with greater predetermined external-debt exposure.

The measurement framework retains information that disappears when newspaper text is reduced to a binary indicator or a single aggregate series. AERI records how prevalent, severe and prominent economic-risk reporting is relative to the overall news flow; the pillars locate that risk across domains; and ARS shows whether the broader configuration is narrow or diffuse. The divergence among these measures is itself useful for surveillance. A sharp rise in AERI with little movement in ARS points to an intense but relatively concentrated episode, whereas increases in both suggest that the risk environment has become broad as well as severe. The pillar indices then show which domains account for that change.\footnote{For surveillance, the joint movement of AERI, ARS and the pillar indices helps distinguish a concentrated episode from one in which several economic domains move together, guiding the choice of follow-up indicators and country-specific analysis.}

The African setting provides a demanding application for this approach. Countries differ sharply in economic structure, external exposure, media depth, language and the availability of conventional high-frequency data. A common construction improves definitional comparability across these environments, while the pillar composition preserves country-specific information.\footnote{A given headline index value can arise from different domain compositions across countries. Country-level interpretation therefore benefits from reading the aggregate score alongside the pillar profile and the associated coverage diagnostics.} The archive-depth diagnostic in Appendix~\ref{app:coverage_diag} finds little within-country association between AERI and retained article volume, although interpretation of any country series still depends on its source mix and economic setting. More broadly, the results show how newspaper archives can organise high-frequency information where conventional data arrive with delays or at low frequency. The contribution is both a new set of African risk series and a transparent framework for separating the level, composition and breadth of reported economic risk.

\subsection{Limitations}\label{sec:limitations}
The first limitation concerns construct scope. AERI, the pillar indices and ARS describe economic-risk content in the retained news archive. Their levels measure the prevalence, severity and composition of reported economic risk within that information environment. Cross-pillar comparisons likewise describe which concerns are most prominent in the corpus; a high financial-sector pillar mean, for example, indicates greater prominence of financial-sector concerns and reflects both underlying conditions and reporting practices.

A second limitation arises from the newspaper environment itself. Publication volume, editorial priorities, source continuity and language coverage differ across countries and over time. The archive-depth diagnostic finds little within-country association between AERI and log article volume in the eligible sample, which is reassuring about a simple volume channel but leaves source representativeness unresolved. The production dictionaries explicitly cover English, French and Portuguese; reporting in other languages may therefore be captured less completely. Exact-text deduplication removes identical reprints, while edited versions, translations and related stories remain distinct records. Its reported sensitivity results apply to the descriptive series; the PVAR and external-shock specifications continue to use the original benchmark vintage.

The PVAR results carry the usual limitations of recursively identified panel systems. The coefficients summarise average conditional dynamics for the selected countries, and the contemporaneous ordering is imposed by the researcher. The main sign pattern persists under both forward and feedback orderings. These estimates therefore provide descriptive evidence on dynamic relationships under recursive identification, while structural causal identification is addressed separately through the external-shock design.

The external monetary-policy design strengthens identification of the shock, while the exposure measure remains observational. Its cross-section contains 29 African economies, and 2014 external-debt exposure may covary with exchange-rate regimes, trade openness, institutional quality and other persistent characteristics that shape the transmission of global tightening. Country and year-month fixed effects, alternative exposure definitions, leave-one-country-out estimates and falsification tests address several specific concerns, yet some cross-sectional confounding may remain. We therefore use the exercise to validate AERI's economic content and exposure sensitivity, reserving causal interpretation of the debt channel for designs with richer predetermined controls or longer panels.
\clearpage
\phantomsection\addcontentsline{toc}{section}{References}
\begingroup
\small
\singlespacing
\printbibliography[heading=none]
\endgroup
\clearpage\appendix
\counterwithin{equation}{section}
\counterwithin{table}{section}
\counterwithin{figure}{section}
\section{Data Collection Architecture}\label{app:data_collection}
The collection and measurement sequence is: country-site registry; monthly retrieval; text normalisation; relevance classification; severity and textual emphasis weighting; country-month aggregation; pillar diagnostics; and ARS construction. A failed source month is treated as a gap in the observation system, while a retained month with no matched economic content is recorded as a measured zero. The replication extension adds exact-text comparison within country-months, including across outlets, and retains duplicate counts, representative examples and file fingerprints. This step addresses identical article bodies; broader questions of story independence and media representativeness remain part of the measurement limitations discussed in the main text.

The index vintage spans 131 or 132 reported months per country within 2015--2025. Annual country values are available from the authors upon request. Descriptive comparisons generally require at least 100 articles per month. The PVAR sample is defined separately: it starts from the 36-country correlation-screened set, retains country-months with complete model inputs and the required lag structure, and requires at least 24 usable observations per country. The descriptive 100-article threshold is not imposed on PVAR estimation. Formal PCA derivations, detailed pillar descriptions and expanded robustness results are reported in the online supplement.

\section{Archive-Depth Diagnostic}\label{app:coverage_diag}
Table~\ref{tab:archive_depth} examines whether AERI moves mechanically with the number of retained articles after persistent country differences are removed. The first specification includes country fixed effects; the second adds year-month fixed effects. Both use country-clustered standard errors and the 6,488 country-months that satisfy the principal coverage threshold.
\begin{table}[htbp]\centering
\caption{AERI and within-country news-archive depth}\label{tab:archive_depth}
\begin{singlespace}\small
\begin{tabular}{lrrrr}
\toprule
Specification & Coefficient & Std. error & $p$-value & Observations\\
\midrule
Country fixed effects & 1.141 & 2.978 & 0.702 & 6,488\\
Country and year-month fixed effects & -1.474 & 3.153 & 0.640 & 6,488\\
\bottomrule
\end{tabular}
\vspace{2mm}
\parbox{0.92\textwidth}{\footnotesize\textit{Notes:} The dependent variable is AERI and the regressor is the natural logarithm of retained article records in the country-month. The sample contains the 6,488 country-months with at least 100 retained articles, covering 53 countries; Guinea-Bissau has no month meeting the threshold in this vintage. Standard errors are clustered by country. The within-country correlation between AERI and log article volume is 0.035. The second specification also absorbs year-month fixed effects.}
\end{singlespace}
\end{table}


\section{External-Identification Diagnostics}\label{app:external_id}
Table~\ref{tab:external_id_diagnostics} reports falsification and small-cluster inference checks for the externally identified monetary-policy exposure design. The joint future-shock placebo includes the contemporaneous shock--exposure interaction as a control. Pre-shock outcome placebos are estimated without lagged outcome controls that would mechanically overlap the placebo dependent variable. Complete horizon paths, robustness results, leave-one-country-out estimates, finite-cluster calculations, wild-cluster-bootstrap outputs and the underlying monthly shock and exposure files are available from the authors on request and are summarised further in the online supplement.
\input{../05_Tables/Table_C1}

\section{PVAR Sample and Estimation Details}\label{app:pvar}
\subsection{Macroeconomic Data and Pooled Panel Estimation}\label{sec:pooled_method}
The empirical assessment relates monthly news-based risk to headline inflation, exchange-rate depreciation and monthly GDP growth. Its purpose is to test whether the measured domains display economically interpretable dynamics. ARS calibration is completed independently of these macroeconomic outcomes, preserving a clean separation between index construction and subsequent validation.

\subsubsection{Monthly variables and sample construction}
Headline inflation, $\pi_{ct}$, is the supplied monthly observation of the inflation rate. The exchange rate, $e_{ct}$, is expressed as local-currency units per US dollar. Monthly GDP, $G_{ct}$, is the supplied purchasing-power-parity GDP series. Depreciation and monthly GDP growth are defined as
\begin{equation}
 d_{ct}=100[\log e_{ct}-\log e_{c,t-1}],\qquad
 g_{ct}=100[\log G_{ct}-\log G_{c,t-1}].\label{eq:macro_growth}
\end{equation}
Thus a positive $d_{ct}$ denotes depreciation, whereas a positive $g_{ct}$ denotes monthly GDP growth. Headline inflation retains the definition supplied in the macroeconomic panel, while depreciation and GDP growth are the log-difference transformations above. The sample-specific winsorisation bounds are reported below.

The main validation specification uses one news-based risk measure at a time---the inflation-risk pillar, the FX and external-risk pillar, AERI or ARS---together with inflation, depreciation and monthly GDP growth. The estimation sample is distinct from the descriptive sample defined by the 100-article threshold. It begins with the 36-country union selected by the prior correlation screen, then retains country-months with the risk measure, the required macroeconomic inputs and the full lag structure, subject to at least 24 usable observations per country.\footnote{The estimation panel winsorises the supplied macroeconomic transformations at pooled first and ninety-ninth percentiles. In the GDP-augmented sample the retained bounds are $[-2.565,50.934]$ for inflation, $[-5.736,14.926]$ for depreciation and $[-0.819,1.506]$ for monthly GDP growth.} The common GDP-augmented panel contains 4,286 observations for 36 countries and runs from May 2015 to December 2025. The 100-article rule continues to govern the principal descriptive country comparisons and ARS calibration; it is not imposed on the PVAR estimation panel. This distinction explains why a country can appear in the PVAR sample even when few or no months satisfy the descriptive coverage threshold. Using the direct pillar indices in this specification keeps the economic interpretation aligned with the measurement framework in the main text. Supporting transformations and country-level estimation diagnostics are available from the authors on request.

Risk innovations are put on a comparable within-country scale. Let $\bar r_c$ be the country mean in the supplied scaling panel and $s_{r,w}$ the sample standard deviation of $r_{ct}-\bar r_c$ pooled across that panel. We use $\widetilde r_{ct}=r_{ct}/s_{r,w}$; country effects subsequently absorb level differences. A one-unit risk innovation therefore represents one pooled within-country standard deviation, rather than one raw index point. Scaling parameters are fixed for the reported estimates and bootstrap exercises.

\subsubsection{The pooled dynamic system}
For $m$ endogenous variables, the maintained model is
\begin{equation}
 y_{ct}=\alpha_c+\lambda_t+\sum_{\ell=1}^{p}B_\ell y_{c,t-\ell}+\varepsilon_{ct},\qquad
 E(\varepsilon_{ct}\varepsilon_{ct}')=\Sigma.\label{eq:pvar}
\end{equation}
Here $\alpha_c$ captures persistent country differences, and $\lambda_t$ is a full year-month effect that absorbs shocks common to countries in a given month. The year-month effects therefore absorb common calendar-time variation rather than recurring month-of-year seasonality alone. The slope matrices $B_\ell$ are pooled across countries. In the three-variable system, $y_{ct}=(\widetilde r_{ct},\pi_{ct},d_{ct})'$. In the GDP-augmented system, $y_{ct}=(\widetilde r_{ct},\pi_{ct},d_{ct},g_{ct})'$. In particular, the inflation equation in the latter system is
\begin{align}
 \pi_{ct}={}&\alpha_c^\pi+\lambda_t^\pi+
 \sum_{\ell=1}^{p}\bigl(b_{\pi r,\ell}\widetilde r_{c,t-\ell}
 +b_{\pi\pi,\ell}\pi_{c,t-\ell}\notag\\
 &\hspace{35mm}+b_{\pi d,\ell}d_{c,t-\ell}
 +b_{\pi g,\ell}g_{c,t-\ell}\bigr)+\varepsilon_{ct}^\pi.\label{eq:inflation_equation}
\end{align}
The risk, depreciation and GDP equations have the same lag structure with equation-specific coefficients. Each outcome can therefore respond to its own history and to the histories of the other variables. Pooling provides a common dynamic description conditional on fixed effects. The multivariate panel structure follows the empirical framework in \citet{love2006} and \citet{abrigo2016}, implemented here with within-transformed pooled least squares.

Let $X$ stack all $mp$ lagged regressors and $Y$ stack the contemporaneous outcomes on the common estimation sample. Alternating country and year-month projections yield $\ddot X$ and $\ddot Y$. The pooled estimator and residual covariance are
\begin{equation}
 \widehat B=(\ddot X'\ddot X)^+\ddot X'\ddot Y,\quad
 \widehat E=\ddot Y-\ddot X\widehat B,\quad
 \widehat\Sigma=\frac{\widehat E'\widehat E}{n-mp}.\label{eq:pooled_ols}
\end{equation}
The superscript $+$ denotes the Moore--Penrose inverse; the covariance denominator is the implemented lag-regressor adjustment. A common maximum-lag sample is used when comparing candidate orders, preventing changes in the available observations from driving lag choice. Among stable candidates with $p\in\{1,2,3\}$, the criterion is
\begin{equation}
 \mathrm{BIC}(p)=\log|\widehat\Sigma_p|+\frac{\log n}{n}m^2p.\label{eq:bic}
\end{equation}
Stability requires the maximum absolute companion-matrix eigenvalue to lie below the numerical threshold of 0.999. All reported specifications select three lags. Table~\ref{tab:model_diagnostics} reports observations, country counts, roots and accepted bootstrap draws. Fixed-effects estimation with lagged outcomes can retain finite-time bias \citep{nickell1981}; the long monthly dimension reduces some finite-sample concerns but leaves this source of bias relevant to interpretation.

\subsubsection{Identification and impulse responses}
Let $\widehat\Sigma=LL'$ with $L$ lower triangular. Forward systems order risk first, followed by inflation, depreciation and monthly GDP growth where included. This permits a risk innovation to affect the macroeconomic variables within the month, while assigning contemporaneous feedback into risk to subsequent periods. For reverse responses, the order is inflation, depreciation, monthly GDP growth and AERI. The two orderings answer complementary recursive questions about forward risk propagation and macroeconomic feedback.

The moving-average matrices and normalised responses are
\begin{align}
 \Phi_0&=I_m,\qquad
 \Phi_h=\sum_{\ell=1}^{\min(p,h)}\widehat B_\ell\Phi_{h-\ell},\label{eq:ma}\\
 \theta_{kj}(h)&=\frac{e_k'\Phi_hLe_j}{e_j'Le_j},\qquad h=0,\ldots,12,\label{eq:irf}
\end{align}
where $e_j$ is the $j$th coordinate vector, distinct from the exchange-rate notation in equation~\eqref{eq:macro_growth}. The denominator normalises the innovation to a one-unit contemporaneous increase in the selected variable. For forward risk innovations, that unit is one within-country standard deviation. For reverse macroeconomic innovations, it is one percentage point. This convention makes the reported response magnitudes explicit while preserving the natural units used for the macroeconomic innovations.

\subsubsection{Inference and interpretation}
Uncertainty bands are constructed by resampling countries with replacement, retaining each selected country's complete time history, and re-estimating the fixed-effects system. The lag order, measurement transformations and fixed ARS calibration remain unchanged across bootstrap replications. Four hundred accepted stable replications provide the fifth and ninety-fifth percentiles at each horizon. These are pointwise 90 percent intervals, so each band describes uncertainty at a given horizon. Simultaneous inference for the full response path or for the location of its population peak would require a different procedure.

Joint lag tests complement the IRFs. For example, the inflation equation tests
\begin{equation}
 H_0:b_{\pi r,1}=\cdots=b_{\pi r,p}=0,\qquad
 W=(R\widehat\beta)'[R\widehat V R']^+(R\widehat\beta),\label{eq:wald}
\end{equation}
using country-cluster covariance $\widehat V$ and a $\chi_p^2$ reference distribution. The exact covariance calculation is given in Appendix~\ref{app:pvar}. Rejection indicates incremental information in the included lagged variable conditional on the other lags and fixed effects. The recursive response additionally reflects contemporaneous residual covariance and propagation through the full system, so the tables report both forms of evidence.


The 36-country GDP extension is based on the union of countries meeting the pre-specified inflation or exchange-rate correlation screen and contains 4,286 common-sample observations.\footnote{The panel comprises Algeria, Angola, Botswana, Burkina Faso, Burundi, Cameroon, Cape Verde, Comoros, Democratic Republic of Congo, Egypt, Ethiopia, Gambia, Ghana, Guinea, Guinea-Bissau, Ivory Coast, Lesotho, Liberia, Libya, Malawi, Mali, Mauritius, Morocco, Mozambique, Namibia, Nigeria, Sao Tome and Principe, Senegal, Sierra Leone, Somalia, South Africa, Tanzania, Tunisia, Uganda, Zambia and Zimbabwe.} Selection is part of the design, and inference is conditional on it. Table~\ref{tab:model_diagnostics} reports the sample sizes, lag choices, stability roots and accepted bootstrap draws. The GDP-augmented roots lie between 0.919 and 0.923.
\input{../05_Tables/Table_D1}

\section{Supplementary Country Outputs}\label{app:supp_country}
Appendix Table~\ref{tab:country_aeri} reports monthly AERI summary statistics for the illustrative ten-country group used in the annual comparison. The online supplement contains the full 54-country AERI/ARS summary, the two multi-country monthly path figures, detailed PVAR estimation material, the expanded exact-text deduplication results and additional external-identification robustness output.
\input{../05_Tables/Table_E1}

\end{document}